\chapter*{Abstract}

\section*{\tituloPortadaEngVal}
 
This study focuses on creating a tool that, using large language models, is capable of generating behavior trees (BTs) for non-player characters (NPCs) to serve as a foundation for reducing the workload of designers.

To this end, the Unreal Engine 5 video game engine was chosen to implement a series of actions and conditionals that will serve as nodes for the generated BTs, as well as to design a series of scenarios that will serve as test environments in which the BTs can be executed.
Once a set of nodes and execution environments was in place, the tool was designed.
Finally, an implementation was chosen in which the “Mistral-large-3” model processed the behavior written in natural language into pseudocode that was translated into a tree structure in JSON format, while a second, smaller model, the “llama3:8B” model, was dedicated to assigning variables in the Blackboard (the BT’s memory module).
Once the BT is generated, two methods for automatic validation and potential correction have been implemented: a non-deterministic finite automaton (NFA) that validates whether the order of the nodes respects any dependencies between them, and a validator that, upon completion of execution, verifies whether certain objectives set for the behavior have been met.
Once the system was set up, user tests were designed to verify whether the tool is intuitive and generates BTs of a quality similar to those created by a human in less time.
The tests were conducted with a sample of N = 3 users, so no definitive conclusions can be drawn.
However, it was observed that users found the tool easy to use (with an SUS score of 90.83/100) and that it generated BTs in less time than humans, albeit with slightly lower quality, thus meeting the tool’s objective.

\section*{Keywords}

\noindent AI, Behavior Trees, NPC Behavior, Natural Language Processing, LLM 



